\documentclass[11pt]{article}

\usepackage[T1]{fontenc}
\usepackage{lmodern}
\usepackage[a4paper,margin=30mm]{geometry}
\usepackage{amsmath,amssymb,amsthm,mathtools}
\usepackage{microtype}
\usepackage{enumitem}
\usepackage{xcolor}
\usepackage[colorlinks=true,linkcolor=blue!45!black,citecolor=blue!45!black,
  urlcolor=blue!55!black]{hyperref}

\newtheorem{theorem}{Theorem}[section]
\newtheorem{proposition}[theorem]{Proposition}
\newtheorem{lemma}[theorem]{Lemma}
\newtheorem{corollary}[theorem]{Corollary}
\newtheorem{definition}[theorem]{Definition}
\newtheorem{remark}[theorem]{Remark}

\newcommand{\C}{\mathbb C}
\newcommand{\R}{\mathbb R}
\newcommand{\N}{\mathbb N}
\newcommand{\HH}{\mathbb H}
\newcommand{\cB}{\mathcal B}
\newcommand{\cG}{\mathcal G}
\newcommand{\cM}{\mathcal M}
\newcommand{\cU}{\mathcal U}
\newcommand{\im}{\operatorname{Im}}
\newcommand{\D}{\partial}

\title{Residual Differential Independence of Countable Translate Families\\
of Positive Discrete Cauchy Transforms}
\author{}
\date{}

\begin{document}

\maketitle

\begin{abstract}
Let $(\lambda_n)$ be any sequence of distinct real numbers escaping every compact set, let
$(b_n)$ be any positive summable sequence, and let $H\subset\C$ be any prescribed countable set
of distinct translations.  For $u\in\ell^\infty(\N,\R)$ put
\[
  f_u(z)=\sum_{n\geq1}\frac{b_ne^{u_n}}{\lambda_n-z}.
\]
We prove that the set of gauges $u$ for which the family
$(f_u(z+h))_{h\in H}$ is differentially algebraically independent over $\C(z)$ is dense
residual.  Every member of the family of sources is a meromorphic Herglotz--Nevanlinna function.
The proof combines finite real controllability of Cauchy jets with a coefficient-uniform
determinant and a Baire-category argument.  We give the local openness and meromorphic
identity steps explicitly.  The result concerns generic freedom inside a positive Cauchy class;
it is distinct from translation hypercyclicity and from hypertranscendence forced by a fixed
linear difference equation.
\end{abstract}

\section{Introduction}

The translates of a meromorphic function may be constrained for two very different reasons.
They may satisfy a difference equation, as in the classical identity
$\Gamma(z+1)=z\Gamma(z)$, or they may be tied together by an algebraic differential relation
without satisfying an evident difference equation.  Parameterized difference Galois theory gives
strong hypertranscendence results in the first regime; see
Adamczewski, Dreyfus, and Hardouin~\cite{ADH2021}.  Algebraic independence for functions governed
by independent difference operators is developed further by Adamczewski, Dreyfus, Hardouin, and
Wibmer~\cite{ADHW2024}.

The present paper studies a complementary generic regime.  We prescribe a countable translation
set in advance and vary only the positive masses of a discrete Cauchy transform.  The conclusion
is not merely that the source is hypertranscendental: all derivatives of all prescribed distinct
translates are algebraically independent over $\C(z)$.

Translation universality provides a different comparison.  Birkhoff's theorem and its common
hypercyclic variants construct entire functions whose selected translation orbits are dense in
spaces of holomorphic functions~\cite{Birkhoff1929,Tsirivas2014}.  Orbit density does not imply
differential algebraic independence, nor conversely.  Our argument uses Baire category, as many
universality arguments do, but the dense open conditions are finite differential-algebraic rank
conditions rather than approximation conditions.

Here is the main statement.  A sequence $(\lambda_n)_{n\geq1}\subset\R$ is called
\emph{discrete} if its terms are distinct and $|\lambda_n|\to\infty$.

\begin{theorem}[Residual translate-freeness law]\label{thm:main}
Let $(\lambda_n)_{n\geq1}$ be a discrete real sequence, let $b_n>0$ with
$\sum_n b_n<\infty$, and let $H\subset\C$ be countable.  Define
\begin{equation}\label{eq:source}
  f_u(z)=\sum_{n\geq1}\frac{b_ne^{u_n}}{\lambda_n-z},
  \qquad u=(u_n)\in\cB:=\ell^\infty(\N,\R).
\end{equation}
There is a dense residual set $\cG_H\subset\cB$ such that, for every $u\in\cG_H$, the family
\begin{equation}\label{eq:translate-family}
  \bigl(f_u(z+h)\bigr)_{h\in H}
\end{equation}
is differentially algebraically independent over $\C(z)$ in the differential field of
meromorphic functions on $\C$.

Moreover, every $f_u$ is a Herglotz--Nevanlinna function:
\begin{equation}\label{eq:herglotz}
  \im f_u(z)>0\qquad(z\in\HH:=\{z:\im z>0\}).
\end{equation}
\end{theorem}

Thus a prescribed countable translation geometry and positivity of the representing measure are
compatible with complete differential-algebraic freedom.  The theorem is uniform in the support:
no spacing, arithmetic structure, or one-sidedness of $(\lambda_n)$ is used.  In particular,
$\lambda_n=\log p_n$ is one possible specialization, but primes play no role in the proof or the
conclusion.

The argument has three steps.  First, finitely many positive masses control an arbitrary finite
packet of Cauchy jets in all real directions.  Second, one determinant excludes, uniformly in
its complex coefficients, every relation having a fixed finite monomial pattern and coefficient
degree.  Third, the set of such patterns is countable, so Baire's theorem supplies one gauge
escaping all of them.

\section{Differential independence and the Cauchy class}

Let $\cM(\C)$ be the differential field of meromorphic functions on $\C$, with derivation
$\D=d/dz$.  A family $(g_i)_{i\in I}\subset\cM(\C)$ is \emph{differentially algebraically
independent over $\C(z)$} if the jet family
\[
  \{\D^rg_i:i\in I,\ r\geq0\}
\]
is algebraically independent over $\C(z)$.  We use the standard terminology of differential
algebra~\cite{Kolchin1973}.

For $u\in\cB$, the coefficients $a_n(u):=b_ne^{u_n}$ are positive and summable.  Since the
support is discrete, the series in \eqref{eq:source}, together with all termwise derivatives,
converges locally uniformly away from $(\lambda_n)$.  Hence $f_u\in\cM(\C)$ and
\begin{equation}\label{eq:jets}
  \D^rf_u(z)=r!\sum_{n\geq1}\frac{a_n(u)}{(\lambda_n-z)^{r+1}}.
\end{equation}
For $z\in\HH$,
\[
  \im\frac1{\lambda_n-z}=
  \frac{\im z}{|\lambda_n-z|^2}>0.
\]
Summing proves \eqref{eq:herglotz}.  The rest of the paper proves the residual assertion.

\section{Finite real controllability of Cauchy jets}

The positivity constraint means that admissible perturbations of the masses are real.  The next
lemma shows that this causes no loss of finite jet directions.

\begin{lemma}[Real Cauchy-jet spanning]\label{lem:spanning}
Let $\zeta_1,\ldots,\zeta_m\in\HH$ be distinct, and let $r_1,\ldots,r_m$ be positive integers.
Then the vectors
\begin{equation}\label{eq:vectors}
  v(\lambda_n)=
  \bigl((\lambda_n-\zeta_j)^{-k}:
  1\leq j\leq m,\ 1\leq k\leq r_j\bigr)
\end{equation}
span $\prod_{j=1}^m\C^{r_j}$ as a real vector space.  A finite subfamily already spans.
\end{lemma}

\begin{proof}
Suppose that a real linear functional $\ell$ on $\prod_j\C^{r_j}$ annihilates every vector in
\eqref{eq:vectors}.  Every real linear functional on a finite-dimensional complex vector space
has the form $x\mapsto\operatorname{Re}\sum c_{j,k}x_{j,k}$.  Therefore the rational function
\[
  R(t)=\sum_{j=1}^m\sum_{k=1}^{r_j}\frac{c_{j,k}}{(t-\zeta_j)^k}
\]
satisfies $\operatorname{Re}R(\lambda_n)=0$ for every $n$.  Set
\[
  Q(t)=R(t)+\overline{R(\overline t)}.
\]
Then $Q$ is rational and has the infinitely many distinct zeros $\lambda_n$, so $Q\equiv0$.
The first summand has all its poles in $\HH$, while the second has all its poles in the lower
half-plane.  The principal part of $Q$ at every $\zeta_j$ is consequently the principal part of
$R$ there, and must vanish.  Thus all $c_{j,k}$ vanish, so $\ell=0$.  The real span is the whole
space.  Finite dimensionality yields a finite spanning subfamily.
\end{proof}

We record explicitly the local-openness consequence needed below.

\begin{lemma}[Positive-mass local openness]\label{lem:openness}
Fix finitely many distinct nodes $\zeta_j\in\HH$ and finitely many jet orders at those nodes.
For every $u\in\cB$ and every $\varepsilon>0$, varying finitely many coordinates of $u$ inside
the ball $\|u'-u\|_\infty<\varepsilon$ makes the resulting finite jet packet of $f_{u'}$ contain
an open neighbourhood, in its full underlying real vector space, of the packet of $f_u$.
\end{lemma}

\begin{proof}
By Lemma~\ref{lem:spanning}, choose a finite index set $E\subset\N$ such that the mass-to-jet map
\[
  A:\R^E\longrightarrow\prod_j\C^{r_j},\qquad
  (\delta a_n)_{n\in E}\longmapsto
  \left(k!\sum_{n\in E}\frac{\delta a_n}
  {(\lambda_n-\zeta_j)^{k+1}}\right)_{j,k}
\]
is surjective as a real linear map.  A surjective linear map between finite-dimensional real
vector spaces is open.  For $n\in E$, the change of variable
\[
  \theta_n\longmapsto
  \delta a_n=a_n(u)(e^{\theta_n}-1)
\]
is a diffeomorphism from a neighbourhood of $0$ onto a neighbourhood of $0$.  Hence the box
$|\theta_n|<\varepsilon$ maps, after composition with $A$, onto a set containing a neighbourhood
of the origin.  Taking $u'_n=u_n+\theta_n$ on $E$ and $u'_n=u_n$ off $E$ proves the claim.  All
perturbed masses remain positive.
\end{proof}

\section{Coefficient-uniform escape}

A single differential relation contains only finitely many translates and jets.  Fix distinct
$h_1,\ldots,h_m\in H$, nonnegative jet bounds $R_1,\ldots,R_m$, and distinct monomials
$M_1,\ldots,M_s$ in the variables
\[
  X_{i,j},\qquad 1\leq i\leq m,quad 0\leq j\leq R_i.
\]
Also fix a nonnegative integer $q$ and put $N=s(q+1)$.  This data will be called a
\emph{finite relation pattern}.

\begin{lemma}[Finite polynomial evaluation]\label{lem:evaluation}
For every finite relation pattern there are points $z_1,\ldots,z_N\in\HH$ such that
all $z_a+h_i$ lie in $\HH$, all these $Nm$ nodes are distinct, and the polynomial
\begin{equation}\label{eq:detpoly}
  \mathcal D(x_1,\ldots,x_N)=
  \det\left[z_a^dM_r(x_a)\right]_{
    1\leq a\leq N,\ (r,d)\in\{1,\ldots,s\}\times\{0,\ldots,q\}}
\end{equation}
on $N$ copies of the jet space is not identically zero.
\end{lemma}

\begin{proof}
The $N$ polynomials $z^dM_r(x)$ are distinct monomials in $(z,x)$, hence linearly independent.
The standard evaluation lemma for linearly independent functions gives $N$ evaluation pairs
$(z_a,x_a)$ for which their square evaluation matrix is invertible: inductively, if no next point
increased the evaluation rank, a nonzero linear combination of the functions would vanish
everywhere.  We may choose the $z_a$ in any nonempty open upper half-plane and avoid the finitely
many affine equalities $z_a+h_i=z_b+h_j$.  Polynomial independence persists on every nonempty
open subset.  The resulting value of \eqref{eq:detpoly} is nonzero, so the determinant polynomial
is not identically zero.
\end{proof}

For the points supplied by Lemma~\ref{lem:evaluation}, define
\begin{equation}\label{eq:actualdet}
  \Delta(u)=\det\left[
    z_a^dM_r\bigl((\D^jf_u(z_a+h_i))_{i,j}\bigr)
  \right]_{a,(r,d)}.
\end{equation}

\begin{proposition}[Pattern escape]\label{prop:escape}
For every finite relation pattern, the set
\[
  \cU_{\mathrm{pattern}}=\{u\in\cB:\Delta(u)\ne0\}
\]
is open and dense in $\cB$.
\end{proposition}

\begin{proof}
Fix $u\in\cB$ and $\varepsilon>0$.  Apply Lemma~\ref{lem:openness} simultaneously to the finite
set of nodes $z_a+h_i$ and all jet orders occurring in the monomials.  The attainable packets
inside the $\varepsilon$-ball contain a real-open neighbourhood of the current packet.  By
Lemma~\ref{lem:evaluation}, $\mathcal D$ is a nonzero complex polynomial.  Such a polynomial
cannot vanish on a nonempty real-open subset of a complex vector space.  Therefore an attainable
packet has nonzero determinant, proving density.

To prove openness, fix a bounded gauge ball.  At the finitely many upper-half-plane evaluation
nodes, formula \eqref{eq:jets} has a common summable majorant, uniformly for gauges in that ball.
Each required jet depends continuously on $u$ in the uniform norm.  Thus $\Delta$ is continuous,
and its nonzero set is open.
\end{proof}

The determinant treats all complex coefficients with the fixed support and degree bound at once.
Indeed, if
\begin{equation}\label{eq:relation-pattern}
  \sum_{r=1}^s\sum_{d=0}^q c_{r,d}z^d
  M_r\bigl((\D^jf_u(z+h_i))_{i,j}\bigr)=0,
\end{equation}
then evaluation at $z_1,\ldots,z_N$ gives the matrix in \eqref{eq:actualdet} times the coefficient
vector $(c_{r,d})$.

\section{Proof of the residual law}

\begin{proof}[Proof of Theorem~\ref{thm:main}]
Because $H$ is countable, there are only countably many finite relation patterns: a pattern is
specified by finitely many elements of $H$, finitely many integer jet bounds, finitely many
monomials with integer exponent vectors, and one integer coefficient-degree bound.  By
Proposition~\ref{prop:escape}, each pattern supplies an open dense subset of the complete metric
space $\cB$.  Their intersection
\[
  \cG_H=\bigcap_{\mathrm{patterns}}\cU_{\mathrm{pattern}}
\]
is dense residual by the Baire category theorem.

Fix $u\in\cG_H$ and suppose that finitely many functions in \eqref{eq:translate-family} satisfy a
nonzero differential polynomial relation over $\C(z)$.  Clearing the finitely many rational
denominators produces a relation of the form \eqref{eq:relation-pattern} for one of the enumerated
patterns.  Since $u$ lies in the corresponding open set, its evaluation matrix is invertible.
The coefficient vector must therefore be zero, a contradiction.

For completeness, there is no hidden domain issue in this evaluation step.  The series defining
$f_u$ and its translates are globally meromorphic.  If the assumed relation is formulated only
as an identity of germs on a sufficiently high half-plane, its left-hand side is a globally
meromorphic function that vanishes on a nonempty open set.  The meromorphic identity theorem
extends the identity to every regular point, in particular to the fixed evaluation points.
Finally, \eqref{eq:herglotz} was proved in Section~2.
\end{proof}

\begin{corollary}[Finite orbit rank]\label{cor:rank}
Under the hypotheses of Theorem~\ref{thm:main}, for residual $u$ and every finite
$F\subset H$,
\[
  \operatorname{dtrdeg}_{\C(z)}
  \C(z)\langle f_u(z+h):h\in F\rangle=|F|.
\]
In particular, every translate is hypertranscendental over the differential field generated by
any finite set of the other translates.
\end{corollary}

\begin{proof}
Differential algebraic independence gives both assertions directly.
\end{proof}

\begin{corollary}[Rational gauge stability]\label{cor:gauge}
Let $u\in\cG_H$ and choose arbitrary rational functions $r_h\in\C(z)$ for $h\in H$.  Then
$(f_u(z+h)+r_h)_{h\in H}$ is differentially algebraically independent over $\C(z)$ and
\[
  \C(z)\langle f_u(z+h)+r_h:h\in H\rangle
  =\C(z)\langle f_u(z+h):h\in H\rangle
\]
as embedded differential fields.
\end{corollary}

\begin{proof}
Every putative relation involves a finite jet packet.  On its polynomial ring, substitution by
$X_{h,j}\mapsto X_{h,j}+\D^jr_h$ is an automorphism, with inverse obtained by subtraction.  It
therefore preserves differential algebraic independence.  Since all derivatives of each $r_h$
belong to the base field, each displayed field contains the generators of the other.
\end{proof}

\section{Comparison and scope}

Theorem~\ref{thm:main} is a generic existence theorem inside a concrete cone of positive
meromorphic functions.  Its content is the simultaneous compatibility of three requirements:
the support and base masses are fixed in advance, all masses remain positive, and an arbitrary
prescribed countable translate family is differentially algebraically independent.  The local
controllability lemma is the mechanism joining positivity to differential freedom.

This differs from three nearby themes.

\begin{enumerate}[leftmargin=*,itemsep=0.5em]
  \item Common hypercyclicity for translation operators concerns density of selected translation
  orbits in a function space~\cite{Tsirivas2014}.  The present rank conditions concern polynomial
  identities among all derivatives of a fixed countable family.

  \item The results of Adamczewski--Dreyfus--Hardouin concern hypertranscendence of solutions to
  fixed linear shift, $q$-difference, and Mahler equations and use parameterized difference
  Galois theory~\cite{ADH2021}.  Our residual sources are not assumed to solve any difference
  equation; the theorem instead proves that their prescribed orbit coordinates satisfy no
  algebraic differential relation at all.

  \item Adamczewski--Dreyfus--Hardouin--Wibmer prove algebraic independence for transcendental
  solutions governed by suitably independent operators~\cite{ADHW2024}.  Here there may be
  arbitrarily many prescribed translations, with no independence condition on their additive
  parameters beyond distinctness, and the conclusion includes every derivative.
\end{enumerate}

The theorem does not provide a closed-form or canonical gauge: $u$ is selected from a residual
set.  It does not say that every Herglotz function has a free translation family, and it does not
identify arithmetic structure in a special support.  In particular, choosing logarithms of
primes as the support does not create a consequence for the Riemann zeta function.  Any such
application would require a separate theorem tying the residual source to Euler-product or zero
data; no such bridge is asserted here.

Two natural problems remain.  First, determine whether the residual set can be replaced by an
explicitly defined gauge for broad classes of supports.  Second, determine which additional
regularity or asymptotic constraints on the masses are compatible with the same countable
translate-freeness law.  Those questions concern the boundary between generic differential
freedom and source-specific rigidity.

\begin{thebibliography}{99}

\bibitem{ADH2021}
B. Adamczewski, T. Dreyfus, and C. Hardouin,
\emph{Hypertranscendence and linear difference equations},
J. Amer. Math. Soc. \textbf{34} (2021), 475--503.

\bibitem{ADHW2024}
B. Adamczewski, T. Dreyfus, C. Hardouin, and M. Wibmer,
\emph{Algebraic independence and linear difference equations},
J. Eur. Math. Soc. \textbf{26} (2024), 1899--1932.

\bibitem{Birkhoff1929}
G. D. Birkhoff,
\emph{D\'emonstration d'un th\'eor\`eme \'el\'ementaire sur les fonctions enti\`eres},
Comptes Rendus de l'Acad\'emie des Sciences de Paris \textbf{189} (1929), 473--475.

\bibitem{Kolchin1973}
E. R. Kolchin,
\emph{Differential Algebra and Algebraic Groups},
Academic Press, New York, 1973.

\bibitem{Tsirivas2014}
N. Tsirivas,
\emph{Existence of common hypercyclic vectors for translation operators},
arXiv:1411.7815 (2014).

\end{thebibliography}

\end{document}
